\begin{afterword}
\summaryhead

The post opens with Terry Pratchett's witches, who know that the gods exist and so do not believe in them. If dragons were real and zebras unseen, fantasy stories would be about zebras. Readers of fantasy who picture themselves carried into a world of magic, most of whom did not become scientists in a world of science, have no reason to think they would act differently there. Without the capacity to be interested in real things, magic made real would lose its charm, as lottery winners' excitement wears off, probably as soon as spells had to be studied.

The same holds for the Future, which will be another Now; for a commenter excited by Rupert Sheldrake's theory rather than by Special Relativity, which is known to be true; and for New Agers, whose beliefs would lose their point if their rituals worked. The author is not against magic, only against a psychology that would pine for assembly lines in a world of spells. The post ends with a litany: whatever one hopes to do anywhere, one may as well do in reality.

\responsehead

The post's thesis is a conditional, and it is plausible: someone who cannot enjoy real things would not long enjoy magic made real, because it would then be one more real thing. The post's one appeal to research, the lottery winners, points to research that partly supports it. The best-known study found that winners were no happier than controls and took less pleasure in mundane events. The details the post adds, six months and the winners' own expectations, are not in that study's abstract, and a larger study published in 2020 found that winners' life satisfaction stayed higher for more than a decade, though the effects on happiness were smaller.

The problem is how the post applies it. The conditional applies only to people who lack the capacity to enjoy the real. With the fantasy readers the post keeps to the conditional: it asks why readers who did not become scientists would act differently in a world of magic, and answers ``If they don't have the scientific attitude,'' without saying how many readers are like this. With a named commenter and with New Agers it drops the condition and states their motives as fact.

The post states as fact that the commenter would not care about Sheldrake's theory if it were taught in schools, and offers no evidence beyond that interest. The comment gave its own reason: that reductionism could not explain protein folding.

The paragraph on New Agers is the thesis restated as a fact about a community, with no evidence about the community.

In short: a plausible conditional, that someone who cannot enjoy the real would not long enjoy magic made real, which the post keeps as a condition for fantasy readers but applies to a named commenter and to New Agers as fact, with no evidence about their motives.
\end{afterword}
